\chapter{Experimentální uspořádání}

\subsection{Optické měřící přístroje}

\subsubsection{Bruker VERTEX 80v}
Jedná se o je vakuový FTIR spektrofotometr, který obsahuje Michelsonův interferometr, se zrcadlem kmitajícím na vzduchovém polštáři. Během všech měření byl používán speciální nástavec, díky němuž bylo možné provádět správné absolutní měření propustnosti. V průběhu měření jsou vzorky umístěny svisle na držáku v měřící komoře, ve které je po dobu měření udržováno vakuum 2,51\,hPa, což slouží k minimalizace vlivu okolního prostředí a vlhkosti na průběh měření. Výrobcem udávaní spektrání rozlišení spektrometru je minimálně 0,2\,cm$^{-1}$ při běžných laboratorních podmínkách. Spektrální rozsah spektrofotometru je 1333 -- 27000\,cm$^{-1}$, ale při požití vhodných rozšíření je možné měřit i do FIR a UV/VIS oblastí \cite{vertex}.

FIXME: Vyfotit obrázek.

%\begin{figure}
%  \centering
%  \includegraphics[width=120mm]{vertex80v.png}
%  \caption{Spektrofotometr Bruker VERTEX 80v.}
%  \label{verteximg}
%\end{figure}

\subsubsection{LAMBDA 1050 UV/Vis/NIR}
LAMBDA 1050 UV/Vis/NIR od firmy PerkinElmer je dvoukanálový spektrofotometr. Jako zdroj slouží halogen-rtuťová a deuteriová lampa.  Přístroj obsahuje dva monochromátory (holografické difrakční mřížky s 1440 proužky,vrypy? na milimetr pro UV/Vis oblast a s 360 proužky na milimetr pro IR oblast)  (FIXME: holographic grating monochromator with 1440 lines/mm UV/Vis blazed at 240 nm and 360 lines/mm NIR blazed at 1100 nm,). Pro dělení svazku na dva kanály je použit chopper. Udávaný spektrální rozsah je 175--3300\,nm, nicméně pro měření pod 185\,nm je potřeba chlazení dusíkem. Použitý detektor pro viditelnou a UV oblast je fotonásobič R6872. Pro oblast 860--1800\,nm je použit InGaAs detektor, chlazený peltierovým článkem a pro oblast 1800--2500\,nm PbS detektor, taktéž chlazený peltierovým článkem. Spektrální rozlišení výrobce udává jako minimálně 0,05\,nm. Přístroj obsahuje optickou kompenzaci pro různé tloušťky vzorků \cite{lambda}.

FIXME: Vyfotit obrázek.

%\begin{figure}
%  \centering
%  \includegraphics[width=120mm]{LAMBDA.jpg}
%  \caption{Spektrofotometr LAMBDA 45 UV/Vis.}
%  \label{lambdaimg}
%\end{figure}

\subsubsection{Jobin-Yvon UVISEL}
Schéma elipsometru Jobin-Yvon UVISEL je znázorněno na schématu \ref{fig:elipsometrimg} (FIXME: udělat schéma v inkscapu). Udávaný spektrální rozsah tohoto přístroje je od 190\,nm do 2000\,nm. Hlavy elipsometru jsou umístěny na goniometru, z toho důvodu může být měření prováděno pod různými úhly. Běžně se měří úhly dopadu 55$^\circ$, 60$^\circ$, 65$^\circ$, 70$^\circ$ a 75$^\circ$. Jedná se o fázově modulovaný elipsometr. Měření probíhá tak, že z lampy L vychází chromatické nepolarizované světlo, které poté prochází přes polarizátor P. Tam dojde k jeho lineární polarizaci, odráží se od vzorku a přes kompenzátor C, analyzátor A a monochromátor dopadá na detektor D.

FIXME: Taky možná fotku?

%\begin{figure}
%  \centering
%  \includegraphics[width=120mm]{schema-elipsometru.png}
%  \caption{Schéma elipsometru UVISEL (P--polarizátor, C--kompenzátor, D--detektor, L--lampa, M--modulátor, A--analyzátor, S--substrát), obrázek převzat z \cite{fialova2009}.}
%  \label{fig:elipsometrimg}
%\end{figure}

\subsection{Depoziční aparatura}
FIXME:tohle je v podstatě opsáno z bakalářky... zkontrolovat jestli to všechni stále sedí!

Všechny vrstvy byly deponovány na Ústavu fyzikální elektroniky Masarykovy University v reaktoru R2. 

Reaktor R2 je válcový ocelový reaktor a má planární elektrody. Jeho vnitřní průměr je 49\,cm a má výšku 24,6\,cm. Horní i dolní elektroda jsou kruhového tvaru. Průměr horní elektrody je 38\,cm, průměr dolní elektrody pak 40\,cm. Horní elektroda je uzemněná. Plyny používané při depozici jsou do reaktoru přiváděny přes otvory v horní elektrodě, kde jsou rozmístěny středově symetricky v kru\-hu o průměru 18\,cm. Deponovaný vzorek je umístěn přímo na dolní elektrodě. Na dolní elektrodu je také přiváděno vysokofrekvenční napětí. Jako zdroj RF napětí byl použit generátor Cesar 133 firmy Dressler pracující na frekvrenci 13,56\,MHz.

Tlak v aparatuře je měřen kapacitronem Leybold. Průtoky plynů do reaktoru jsou regulovány průtokoměry Hastlings. 

\subsection{Depoziční podmínky}

Nezávislé depožiční podmínky pro všechny depozice je tlak $p$, výkon $P$, čas depozice $t$ a průtoky jednotlivých plynů $Q_{\mathrm{CH_4}}$, $Q_{\mathrm{H_2}}$, $Q_{\mathrm{Ar}}$, všechny depoziční podmínky jsoun shrnuty v tabulce. Předpětí $U_b$ pak není nezávislý parametr, ale závisí na výkonu přibližně jako $U_b \sim \sqrt{P}$, nicméně je také zahrnuto v tabulce protože se jedná o velmi důležitý parametr pro růst vrstvy.

\begin{table}[htbp]
 \centering
 \begin{tabular}{lccccccc}
	\hline
	{Vrstva} & {$p$\,[Pa]} & {$P$\,[W]} & {$U_b$\,[V]} & {$t$\,[min]} & {$Q_{\mathrm{CH_4}}$} & {$Q_\mathrm{H_2}$} & {$Q_\mathrm{Ar}$} \\
	& & & & & {[sccm]} & {[sccm]} & {[sccm]} \\
	\hline\hline
	CH30A     & 12 & 100 & 75  & 30   & 8,7 & 5 & 0 \\
	CH30A\_280 & 12 & 100 & 75  & 30   & 8,7 & 5 & 0 \\
	CH30A\_340 & 12 & 100 & 75  & 30   & 8,7 & 5 & 0 \\
	CH83A     & 12 & 75  & 115 & 120  & 8   & 0 & 5 \\
	CH87A     & 8  & 60  & 120 & 120  & 8   & 5 & 0 \\
	CH88A     & 8  & 60  & 140 & 4*30 & 8   & 5 & 0	\\
	CH90A     & 8  & 65  & 145 & 30   & 8   & 5 & 0 \\
	\hline
 \end{tabular}
 \caption{Depoziční podmínky, Q$_\mathrm{{HMDSO}}$ je průtok HMDSO, Q$_\mathrm{O2}$ je průtok kyslíku, $P$ je výkon, $U_b$ je samopředpětí, $p$ je tlak a $t$ je délka depozice}
\end{table}

Všechny vrstvy byly také před depozicí plazmově čištěny v ???? plazmě. Délku a druh plazmkového čištění shrnuje tabulka.....

\cleardoublepage
